\section{Conclusion}
\label{sec:conclusion}

\subsection{Summary of Contributions}
In this paper, we have addressed the critical challenge of data imbalance in intelligent mechanical fault diagnosis by developing the CAC-CycleGAN-WGP framework. Our primary contribution is the integration of a Conditional Auxiliary Classifier (CAC) into a cycle-consistent architecture, which enables high-fidelity cross-domain translation from healthy vibration signals to multiple fault categories. Unlike traditional generative models that synthesize signals from random noise, our approach utilizes the structural information of real normal signals as a baseline, ensuring that the generated fault signatures are physically consistent and contextually relevant. Furthermore, we incorporated the Wasserstein distance with Gradient Penalty (WGP) to stabilize the bidirectional training process, effectively mitigating common issues such as mode collapse and gradient vanishing. To validate the quality of the augmented data, we implemented a hierarchical evaluation scheme using Stacked Autoencoders (SAE) for feature reconstruction analysis and Support Vector Machines (SVM) for diagnostic verification, providing a more rigorous assessment than simple visual inspection of spectra.

\subsection{Key Findings and Future Work}
The experimental validation on two benchmark datasets—the CWRU bearing dataset and the IEEE PHM 2009 gearbox dataset—demonstrates the superior performance of the proposed method. Key findings indicate that the CAC-CycleGAN-WGP consistently outperforms state-of-the-art benchmarks like SMOTE, ADASYN, and standard WGAN-GP across various imbalance ratios (BR). Specifically, in extreme cases (e.g., BR = 0.005), our method improves diagnostic accuracy from under 50\% to over 98\%, proving its effectiveness in low-sample regimes. The ablation study further confirmed that the cycle-consistency loss is the most vital component for maintaining signal integrity, while the CAC is essential for class discriminability. Despite these achievements, future work will focus on extending this framework to handle non-stationary signals under varying speed and load conditions more dynamically. Additionally, we aim to explore the potential of self-supervised learning to further reduce the dependency on initial minority labels, moving toward a truly zero-shot fault generation paradigm for emerging industrial equipment where no fault data has yet been recorded.
